\documentclass[11pt]{article}
\usepackage[a4paper,margin=25mm]{geometry}
\usepackage{amsmath,amssymb,amsthm}
\newtheorem{theorem}{Theorem}
\newtheorem{lemma}{Lemma}
\title{Structural certificates and a correction for the $A_{11}$ atlas table}
\author{L054 finite application}\date{}
\begin{document}\maketitle

\begin{theorem}
Relative to the acquired numerical catalogue and structural registry,
$787$ of the $788$ structural labels in the retained companion table
headed $A_{11}$ hold for their corresponding embedded subgroups.
Row $186$ prints $2\cdot A_5:2$. Under the article's convention of left
association, this expression is impossible in $A_{11}$; the corresponding
subgroup instead has structure $(A_5\times2)\cdot2$, equivalently
$2\cdot(A_5:2)$ with the displayed parentheses.
An extension expression asserts a normal subgroup, quotient and the
stated splitting behaviour; it need not specify a unique isomorphism class.
\end{theorem}

\paragraph{Source and acquired results.}
The companion table by Thomas Connor and Dimitri Leemans is dated
27 November 2012 and accompanies \cite{atlas}. Its numerical catalogue,
source correspondence, embedded generators and orders are acquired from
the separately accepted computation in degree $11$. There are $788$
conjugacy classes and $3304$ weighted maximal-subgroup edges. The subgroup
census and incidence counts are not repeated here. The registry contains
$638$ finite models reconstructed from previously accepted applications;
the files used to reconstruct them are bound to those reviews. The acquired
hypotheses and qualifications remain part of the present statement.

Section 3.2 of the article explicitly gives left association for
unparenthesized extensions and distinguishes split from nonsplit
extensions. We acquire that convention from its separately reviewed native
reading. A new source audit joins the native powers, subscripts,
parentheses and operators of all $788$ labels to the retained expressions.
The table has $41$ pages and $17498$ retained text states. Native pages
$1,2,11,13,18,40,41$ were inspected directly; the remaining pages have
position checks, not a claim of visual inspection. The position audit uses
the same pypdf decoder as the flat extraction. Row $357$ places redundant
parentheses around the single atom $D_{10}$; only that exact glyph sequence
is normalized. The final caption is checked explicitly and excluded from
the last cell. Every printed expression, including row $186$, is retained
unchanged. The claim concerns this pinned edition of the table.

\begin{lemma}
Let $H$ be generated by specified permutations and let $S$ be the union of
their supports, with $n=|S|$. If the order of this same embedded group is
$n!$, then $H=\operatorname{Sym}(S)$. If all its specified generators are
even and its order is $n!/2$, then $H=\operatorname{Alt}(S)$.
\end{lemma}
\begin{proof}
Every generator fixes the complement of $S$, so $H\leq\operatorname{Sym}(S)$.
Equality of finite orders proves the first assertion. For the second,
the sign homomorphism gives $H\leq\operatorname{Alt}(S)$, and equality of
orders again gives equality of groups. The application first joins each
source row to the accepted numerical record and exact generator tuple.
The order used here is consequently the order of that group, rather than
an inference from its printed name.
\end{proof}

This identifies rows $1,2,4,12,39,138,414,706$. Symmetric subgroups embedded
in an alternating group can have extra moved points, so a second criterion
is useful for eight further rows.

\begin{lemma}
Let $H$ act on a finite set $X$, let $S$ be invariant, and put $n=|S|\geq3$.
Suppose that $t$ is an involution on $X\setminus S$, possibly the identity,
and that each specified generator $h$ acts on the complement as
$t^{\epsilon(h|_S)}$, where $\epsilon$ is the additive sign in $\mathbb Z/2$.
If the acquired order of $H$ is $n!$, restriction identifies $H$ with
$\operatorname{Sym}(S)$.
\end{lemma}
\begin{proof}
The sign map is a homomorphism, so the displayed relation on generators
holds on every word and its inverse. Any element acting trivially on $S$
has sign zero and therefore fixes the complement as well. Thus restriction
is injective. Its image in $\operatorname{Sym}(S)$ has order $n!$, so it
is the whole symmetric group. The implementation finds actual generator
orbits, computes each restriction sign by inversion parity, and checks
the same complement involution on every generator. It does not enumerate
$H$ or infer its order from this test.
\end{proof}

The second lemma applies to rows $3,9,26,82,277,287,762,766$. It is the
usual argument for the graph of a homomorphism. The adaptation here
recognizes the sign action on the complementary points.

\begin{lemma}
Let $M$ and $H$ be finite groups given by complete closures of specified
permutation generators. A bijection $f:M\to H$ with $f(1)=1$ and
$f(xg)=f(x)f(g)$ for every $x\in M$ and every specified generator $g$
is an isomorphism.
\end{lemma}
\begin{proof}
Every element is a positive word in the generators: in a finite group
the inverse of a generator is a nonnegative power of it. Induction on
such a word proves $f(xy)=f(x)f(y)$. The reader constructs the proposed
images throughout the complete closure, checks consistency on repeated
elements and checks every right-generator edge. It verifies that the
image is exactly the target subgroup. Thus the checked map is a bijective
homomorphism, rather than only a correspondence between generating sets.
\end{proof}

\paragraph{Products and extensions.}
The new registry has $155$ models: $96$ split extensions, $30$ direct
products, $28$ nonsplit extensions and one dicyclic group. Each model is
proved before any later model or source occurrence can use it. For an
extension model $H$, the reader closes the displayed subgroup $N$,
checks $N\triangleleft H$, and transports an earlier proved model onto
$N$. A split certificate supplies a subgroup $K$ of the indicated type,
checks $N\cap K=1$ and $|N||K|=|H|$, and hence proves $H=NK$.
For the colon notation it also witnesses nontrivial conjugation of $K$
on $N$. A direct product instead checks commutation for all pairs of
factor generators; the factors then commute elementwise.

Fifteen nonsplit models have quotient $C_2$. A complement exists exactly
when there is an involution outside $N$. Indeed, such an element generates
a subgroup disjoint from $N$ whose product with $N$ has order $|H|$.
The reader checks every outside element, so absence of an involution
proves nonsplitting.

The other thirteen nonsplit models have quotient $S_3$. The quotient
certificate identifies generators with defining relations
$a^2=b^3=(ab)^2=1$. A complement provides lifts satisfying the same
relations. Conversely, such lifts define a homomorphism from $S_3$ whose
composition with the quotient map is the identity, and hence a section.
All possible lifts are considered. Those failing the first two relations
may be rejected before the third is checked; the recorded pair count
still includes the full original Cartesian product. Every pair fails
at least one relation, which excludes a section.

For the source symbol $Q_{28}$, the new literal model has generators
$x,y$ of a group of order $28$, with
\[
 |x|=14,\qquad y^2=x^7,\qquad y^{-1}xy=x^{-1}.
\]
These relations reduce every word to $x^i$ or $x^iy$, with $0\leq i<14$.
Thus the presented dicyclic group has at most $28$ elements. The actual
generated permutation group has exactly $28$, so its canonical quotient
map is an isomorphism. The inherited predicate already checked these
relations, complete generation and cardinality. Its finite allowed list
was extended from orders $8,12,16,20$ to include $28$; no relation was
weakened. The earlier prepared requests were not executed.

\paragraph{The separate correction.}
The accepted companion note \cite{correction} proves that a nonsplit
extension $2\cdot A_5$ cannot embed in $A_{11}$. Such an extension is
perfect and has a central involution. Its restriction to the support of
that involution would lie in $C_2\wr S_m$ for $m=2$ or $4$, a solvable
group. A perfect group has trivial image there, contradicting the
involution's nontrivial restriction. The note also identifies the actual
row-$186$ subgroup as the parity pullback
\[
 \{(\sigma,k)\in S_5\times C_4:
                  \operatorname{sgn}(\sigma)=(-1)^k\}.
\]
Its subgroup of index two is $A_5\times C_2$, and every element outside
it has nontrivial square in the $C_4$ coordinate. The extension is
therefore nonsplit. The full proof and its literal $240$-element
certificate are acquired from the distinct review of that note. The
present reader replays the certificate on the exact saved generators.
The correction remains separate from the $787$ recovered printed labels;
it is not an open-conjecture resolution.

\paragraph{Verification and attribution.}
The complete literal check passed in $23.697$ seconds, with $788$ source
joins, $155$ new models, $1051$ transports and $3892103$ generator edges.
It performed $19706$ outside-element tests and considered $87108$ quotient
lift pairs. The support and restriction lemmas each account for eight
rows; $771$ source occurrences use literal transports, and row $186$
uses its separate corrected certificate. Inputs remained unchanged and
stderr was empty. The native source wrapper independently replayed all
$17498$ retained text states and all $788$ glyph joins.

Fourteen damaged source inputs are rejected, including changes to the
printed row-$186$ assertion and the final caption. Five narrow parenthesis
controls permit only the single $D_{10}$ normalization. The two short
recognition lemmas have their separate literal controls. Eleven damaged
$Q_{28}$ certificates are rejected in ordinary and optimized Python,
including a cyclic group of order $28$ satisfying the order and square
conditions but failing the inversion relation.

The numerical catalogue, source expressions and notation belong to
Connor and Leemans. The finite presentation, restriction, transport and
extension arguments are standard group theory or acquired common results.
The work here provides the explicit joins and finite certificates for
this table, together with the separately proved correction. Candidate
searches supply witnesses only; the literal reader establishes their
assertions. Distinct final source admission remains to be recorded.
No Lean formalization, whole-article completion or compression measurement
is claimed here. Other large companion tables remain in scope for L054.

\begin{thebibliography}{2}
\bibitem{atlas} T. Connor and D. Leemans,
An atlas of subgroup lattices of finite almost simple groups,
\emph{Ars Mathematica Contemporanea} \textbf{8} (2015), 259--266.
The retained companion table is headed $A_{11}$ and dated 27 November 2012.
\bibitem{correction} human\_compression project,
\emph{A structural label in the subgroup table of $A_{11}$},
companion proof and literal certificate, acquired in the present packet.
\end{thebibliography}
\end{document}
